\documentclass{report}
\usepackage{amsmath,amssymb}
\setlength{\parindent}{0mm}
\setlength{\parskip}{1em}
\begin{document}
\begin{center}
\rule{6in}{1pt} \
{\large Keith W Kelly \\
{\bf A scalable algorithm for inverse medium problems with multiple sources}}

201 E 24th St \\ Austin \\ TX 78712
\\
{\tt keith@ices.utexas.edu}\\
George Biros\end{center}

We consider the problem of acoustic scattering as described by the
free-space, time-harmonic scalar wave equation given by
\begin{equation}
\Delta v(\mathbf{r}) + k^2(\mathbf{r})v(\mathbf{r}) = -s(\mathbf{r}),
\label{eq:helm}
\end{equation}
along with radiation boundary conditions. Here, $r$ is a point in
$\mathbb{R}^3$, $s$ is the source term, and $k$ is the wavenumber.
Our formulation is based on potential theory. First we write $k$ as $k_0
+ k(r) + \eta(r)$, where $k_0$ is a constant value, $k(r)$ is a known
function of space, and $\eta$ is the unknown medium perturbation.
Then, given the sources $s(\mathbf{r})$ and detectors outside the
perturbation $\eta$ we solve an inverse problem based on (\ref{eq:helm})
for $\eta$.
We observe the total field only at some (finite dimensional set of) sensor positions.
We use multiple events, where we consider independent activations or
different sources at different times.
This problem is common in many areas of science and engineering such as
seismic imaging, subsurface imaging, impedance tomography,
non-destructive evaluation, and diffuse optical tomography.

We describe an algorithm for efficiently solving the inverse scattering
problem for the low-frequency time-harmonic scalar wave equation with
multi-point illumination.
The following approach is based on the FaIMS algorithm described by
Chaillat and Biros in [1].
First, we compress the number of incident fields (computed by solving the
variable coefficient Helmholtz equation with point sources) using a
randomized QR factorization to compute a low rank approximation.
By compressing the incident fields, we greatly reduce the problem size
and by using a randomized QR factorization we can compute an
approximation to within a specified tolerance, ensuring that there is not
significant information loss.
After compressing the incident fields, we can compute the medium
perturbation $\eta$ by solving $\Delta \phi(\mathbf{r}) +
k_0^2\phi(\mathbf{r}) = -\eta(\pmb{r})v(\pmb{r})$, where the scattered
field $\phi$ is measured as the difference between the incident field and
the measured total field.
We transform the Helmholtz equation above into a linear integral equation to obtain
\begin{equation}
\phi(\pmb{r}') = \int_{\mathbb{R}^3}
G(\pmb{r},\pmb{r}')\eta(\pmb{r})(\phi(\pmb{r}) + u(\pmb{r})) d\pmb{r},
\label{eq:ls}
\end{equation}
where $G$ is the Green's function for the Helmholtz operator with wavenumber $k_0$.
Using the Born approximation we can eliminate the dependence of the right
hand side of the Lippmann-Schwinger equation on the scattered field.
In order to solve for $\eta$ we use the conjugate gradient method on the
normal equations with a randomized QR factorization based preconditioner.
To improve the accuracy of the computed solution for slightly larger
perturbations of the background medium, we use an iterated Born
approximation scheme.

We have extended the results of [1] in a few different ways.
In the original paper the background medium was assumed to have a
constant wavenumber except for the perturbation, i.e. $k(r)$ was zero,
but here we allow for variable background medium.
We expand the types of scatterers allowed; the original paper allowed
only point scatterers while we allow for multiple continuous scattering
objects.
We also make use use of the iterated Born approximation in our solver.
Finally, we achieve superior scalability to the approached referenced in the FaIMS paper.
We rapidly evaluate the integral for the forward scatterer using a fast
multipole method for volume potentials in parallel.
Furthermore, we use \texttt{PETSc} for in our implementation for its fast
iterative solvers.

\section*{References}
[1] St\'{e}phanie Chaillat and George Biros. FaIMS: A fast algorithm for
the inverse medium problem with multiple frequencies and multiple sources
for the scalar helmholtz equation. \textit{Journal of Computational
Physics}, 231(20):4403�4421, 2012.


\end{document}
